% !TEX root = ../main.tex

\section{Personalized Federated Learning Problem}
In this section, we introduce the personalized federated learning problem that aims to collaboratively train personalized models for a set of clients using the non-IID private data of all clients in a privacy-preserving manner~\cite{kairouz2019advances, zhao2018federated}.


\nop{
In this section, we introduce the personalized federated learning problem
that involves a set of clients and a central cloud server.
Each client owns a local model and a set of private training data, where the training data of different clients are sampled from distinct distributions. 

Each client owns a personalized local model and a set of private training data collected from distinct distributions. 
The cloud is a central server that coordinates the training process.


effective collaborations between clients in a .

}

\nop{
, each owning a personalized local model and a set of private training data collected from distinct distributions, attempt to train their personalized models by effectively collaborating with each other in a privacy-preserving manner~\cite{kairouz2019advances, zhao2018federated}. 
}

\nop{ under the coordination of a service provider}

\nop{
, where the personalized local models of all clients are in the same type but have distinct model parameters.
}

\nop{
We will illustrate how to derive the \textbf{personalized cloud model} for each client later in this subsection.
}

\nop{
Inspired by several widely adopted formulations in multi-task learning~ \cite{han2015learning,zhang2017survey}, 
}

\nop{We say $w_i$ is the model of the client $\mathcal{C}_i$ when the context is clear.}

Consider $m$ clients $C_1, \ldots, C_m$ that have the same type of models $\mathcal{M}$ personalized by $m$ different sets of model parameters $\mathbf{w_1}, \ldots, \mathbf{w_m}$, respectively. 
Denote by $\mathcal{M}(\mathbf{w_i})$ and $D_i$ $(1 \leq i \leq m)$ the personalized model and the private training data set of client $C_i$, respectively.
These data sets are non-IID, that is, $D_1, \ldots, D_m$ are uniformly sampled from $m$ distinct distributions $P_1, \ldots, P_m$, respectively.
For each client $C_i$, denote by $\mathcal{V}_i$ the performance of $\mathcal{M}(\mathbf{w_i})$ on the distribution $P_i$.  Denote by $\mathcal{V}^*_i$ the best performance model $\mathcal{M}$ can achieve on $P_i$ by considering all possible parameter sets.

The \emph{personalized federated learning problem} aims to collaboratively use the private training data sets $D_1, \ldots, D_m$ to train the \textit{personalized models} $\mathcal{M}(\mathbf{w_1}), \ldots, \mathcal{M}(\mathbf{w_m})$ such that $\mathcal{V}_1, \ldots, \mathcal{V}_m$ are close to $\mathcal{V}^*_1, \ldots, \mathcal{V}^*_m$, respectively, and no private training data of any clients are exposed to any other clients or any third parties.


\nop{In summary, a personalized federated learning process allows each client $C_i$ to have a personalized model $\mathcal{M}(w_i)$ of its own, such that $\mathcal{M}(w_i)$ achieves a good personalized performance on the corresponding data distribution $P_i$.
}

\nop{Formally, let $\delta$ be a non-negative real number, if 
\begin{equation}\nonumber
    \forall i\in\{1,\ldots, m\}, \mathcal{V}^*_i - \mathcal{V}_i\leq \delta,
\end{equation}
we say the personalized federated learning algorithm has $\delta$-accuracy loss.
}

\nop{
A personalized federated learning problem aims to 
}

\nop{
Specifically, we consider $m$ clients $\{\mathcal{C}_1, \ldots, \mathcal{C}_m\}$ that have the same type of personalized models parameterized by $m$ different sets of $d$-dimensional model parameters $\{w_1, \ldots, w_m\}$, respectively. 
Denote by $\{B_1, \ldots, B_m\}$ the private training data sets of the $m$ clients, respectively. These data sets are non-IID across different clients, that is, $B_1, \ldots, B_m$ are sampled from distinct data distributions.
A straight-forward method is to use each of $B_1, \ldots, B_m$ to separately train the personalized models $\{w^{sep}_1, \ldots, w^{sep}_m\}$, respectively.
A personalized federated learning system is a learning process in which all the clients collaboratively train their own personalized models $\{w^{per}_1, \ldots, w^{per}_m\}$, such that 


Another method, 

A conventional method is to collect all the data together and use $B_{all}=B_1\cup \ldots \cup B_m$ to train a model $\mathcal{M}_{all}$.
}

\nop{
Denote by $w_i\in\R^d$, $i\in\{1,\dots,m\}$, the parameters of the personalized local model of the $i$-th client, and by $F_i:\R^d\rightarrow\R$ the training objective function with respect to the private training data of the $i$-th client.
}

\nop{Intuition of problem formulation.}

%Next, we introduce our formulation of this problem that allows each client to train its own personalized model, and boosts the collaboration effectiveness between clients by an attention-inducing function that iteratively encourages similar clients to collaborate more with each other.

To be concrete, denote by $F_i:\R^d\rightarrow\R$ the training objective function that maps the model parameter set $\mathbf{w_i}\in\mathbb{R}^d$ to a real valued training loss with respect to the private training data $D_i$ of client $C_i$.
We formulate the personalized federated learning problem as 
\begin{equation}
\label{eq:our-form}
\min_{W} \ \left\{\mG(W)\coloneqq \sum_{i=1}^m F_i(\mathbf{w_i}) + \lambda\sum_{i<j}^m A(\|\mathbf{w_i} - \mathbf{w_j}\|^2)\right\}, 
\end{equation}
where $W = [\mathbf{w_1},\dots,\mathbf{w_m}]$ is a $d$-by-$m$ dimensional matrix that collects $\mathbf{w_1},\dots,\mathbf{w_m}$ as its columns and $\lambda>0$ is a regularization parameter.

The first term $\sum_{i=1}^m F_i(\mathbf{w_i})$ in Eq.~\eqref{eq:our-form} is the sum of the training losses of the personalized models of all clients. This term allows each client to separately train its own personalized model using its own private training data.
%
The second term improves the collaboration effectiveness between clients by an attention-inducing function $A(\|\mathbf{w_i} - \mathbf{w_j}\|^2)$ defined as follows.

\nop{
Here, the formulation in \eqref{eq:our-form} is a general form that represents a class of objective functions for different user-specified functions $D$.
The second term of \eqref{eq:our-form} will enforce the attentive collaboration mechanism if $D$ is properly chosen to be an \textbf{attention-inducing function} defined as follows.
}

\begin{defi}\label{defi:att-ind}
$A(\|\mathbf{w_i} - \mathbf{w_j}\|^2)$ is an attention-inducing function of $\mathbf{w_i}$ and $\mathbf{w_j}$ if $A:[0,\infty)\rightarrow\R$ is a non-linear function that satisfies the following properties. 
%(i) $\mA(0) = 0$, (ii) $\mA$ is increasing and concave in $[0,\infty)$  
\begin{enumerate}
\item $A$ is increasing and concave on $[0,\infty)$ and $A(0) = 0$;
\item $A$ is continuously differentiable on $(0,\infty)$; and
\item For the derivative $A^\prime$ of $A$, $\lim_{t\rightarrow 0^+} A^\prime(t)$ is finite.
%\item For the derivative $A^\prime$ of $A$, we have $\lim_{t\rightarrow 0^+} A^\prime(t) = a$ for some $a<\infty$.
\end{enumerate}
\end{defi}
\nop{According to the above conditions (i) and (ii), }
\nop{
Obviously, an attention-inducing function $D(x,y)$ is an increasing function of the Euclidean distance between $x$ and $y$, and it is continuously differentiable at any $x,y\in\R^n$. 
}

\nop{The class of attention-inducing functions is broad.}

The attention-inducing function $A(\|\mathbf{w_i} - \mathbf{w_j}\|^2)$ measures the difference between $\mathbf{w_i}$ and $\mathbf{w_j}$ in a non-linear manner. 
A typical example of $A(\|\mathbf{w_i} - \mathbf{w_j}\|^2)$ is the negative exponential function $A(\|\mathbf{w_i} - \mathbf{w_j}\|^2) = 1 - e^{-\|\mathbf{w_i} - \mathbf{w_j}\|^2/\sigma}$ with a hyperparameter $\sigma$. Another example is the smoothly clipped absolute deviation function~\cite{fan2001variable}. One more example is the minimax concave penalty function~\cite{zhang2010nearly}.
We adopt the widely-used negative exponential function for our method in this paper.

%\chu{mention what we choose in our experiment}

\nop{
2) the tamed square root function $A(\|w_i - w_j\|^2) = \|w_i - w_j\|^2/(2\sigma)$, where $\|w_i - w_j\|^2\in[0,\sigma^2]$ and $\sigma$ is a user given parameter;
}

\nop{
as well as $\mA(\|w_i - w_j\|^2) = \sqrt{\|w_i - w_j\|^2}-\sigma/2$ when $\|w_i - w_j\|^2>\sigma^2$ and $\sigma>0$ is a parameter that can be chosen by the user. 


Also, the smoothly clipped absolute deviation (SCAD) function~\cite{fan2001variable} and the minimax concave penalty (MCP) function~\cite{zhang2010nearly}, which are popular for inducing sparse estimators in high-dimensional statistics, are also examples that satisfy (i) and (ii). 
}

As to be illustrated in the next section, our novel use of the attention-inducing function realizes an attentive message passing mechanism that adaptively facilitates collaborations between clients by iteratively encouraging similar clients to collaborate more with each other. The pairwise collaborations boost the performance in personalized federated learning dramatically.


\nop{
is the key to conducting the attentive message passing mechanism that adaptively discovers the underlying collaboration relationships between clients to achieve outstanding collaboration effectiveness.
}

\nop{
We first introduce how to tackle this problem by a general framework named {FedAMP} in Sec.~\ref{sec:alg}, then we discuss a specific instantiation of {FedAMP} as well as a practical heuristic extension called {\sc HeurFedAMP} in Sec.~\ref{sec:instances} and Sec.~\ref{sec:heuristic}, respectively.
Last, we prove the convergence of {FedAMP} for both convex and non-convex local models in Sec.~\ref{sec:analysis}.
}

\nop{
We consider a typical setting of the \textbf{personalized federated learning problem}, where a set of clients, each owning a personalized model and a set of private data collected from distinct distributions, attempt to train their personalized models by effectively collaborating with each other in a privacy-preserving manner under the coordination of a service provider~\cite{zhao2018federated, kairouz2019advances}. 

For the rest of this section, we first introduce how to tackle the personalized federated learning problem by a general framework of federated attentive message passing in Sec.~\ref{sec:alg}, then we discuss a specific instantiation of {FedAMP} as well as a practical heuristic extention called {\sc HeurFedAMP} in Sec.~\ref{sec:instances} and Sec.~\ref{sec:heuristic}, respectively.
Last, we prove the convergence of {FedAMP} for both convex and non-convex client models in Sec.~\ref{sec:analysis}.
}
\nop{
In Sec.~\ref{sec:alg}, we introduce federated attentive message passing ({FedAMP}) as a general framework that tackles the personalized federated learning problem by an attentive collaboration ma, then we discuss a practical instantiation of {FedAMP} in Sec.~\ref{sec:instances} using RBF kernel, 
}

\nop{
Due to the non-IID data across the clients, it is desirable to (i) help each client train a personalized model and (ii) request the server to provide 
personalized guidance to each client. Moreover, in order to boost the effectiveness of federated learning under the non-IID setting, it is beneficial to invoke an attentive mechanism in collaboration. In particular, we hope that the service provider encourages each client to collaborate more with those that share much similarity with it, without unveiling the collaboration graph to the clients. In this section, we propose a novel framework named federated attentive message passing ({FedAMP}) that embraces the said properties.
}